% =========================================================================
% English translation (generated by AI using Claude Opus 4.8 (Max effort)).
% This file was translated from Hungarian to English by AI.
% DISCLAIMER: Please review against the original Hungarian .tex files, before relying on it!
% SOURCE: 2.1 Differenciálegyenletek.tex
% =========================================================================
\documentclass[margin=0px]{article}

\usepackage{listings}
\usepackage[utf8]{inputenc}
\usepackage{graphicx}
\usepackage{float}
\usepackage[a4paper, margin=1in]{geometry}
\usepackage{subcaption}
\usepackage{amsthm}
\usepackage{amssymb}
\usepackage{amsmath}
\usepackage{fancyhdr}

\renewcommand{\figurename}{Figure}
\newenvironment{tetel}[1]{\paragraph{#1 \\}}{}
\newcommand{\R}{\mathbb{R}}

\pagestyle{fancy}
\lhead{\it{CS BSc Final Exam Topics}}
\rhead{2.1 Differential equations}

\title{\textbf{{\Large ELTE IK - Computer Science BSc} \vspace{0.2cm} \\ {\huge Final Exam Topics}} \vspace{0.3cm} \\ 2.1 Differential equations}
\author{}
\date{}

\begin{document}
\maketitle

\begin{center}
\fbox{\parbox{0.92\textwidth}{\small \textit{Note: This document was translated from Hungarian to English by AI. Please verify the information against the original Hungarian source before relying on it.}}}
\end{center}

\section{The initial value problem}
\begin{description}
    \item[Differential equation] \hfill \\
        $ 0 < n \in \mathbb{N}, \ I \subset \R$ open interval, \\
        $ \Omega := I_1 \times ... \times I_n \subset \R^n$, where $ I_1,...,I_n \subset \R$ open interval \\
        $f:I\times\Omega \rightarrow \R^n, \ f \in C $

        Determine the function $ \varphi \in I \rightarrow \Omega$ such that:
        \begin{itemize}
            \item $ D_{\varphi} $ open interval
            \item $ \varphi \in D $
            \item $ \varphi'(x) = f(x, \varphi(x)) \quad (x \in D_{\varphi}) $
        \end{itemize}

        We call this task a differential equation.
    \item[Initial value problem] \hfill \\
        If, in addition to the above, the following are also given: $ \tau \in I$, and $ \xi \in \Omega$ \\
        And the function $\varphi$ additionally satisfies:
        \begin{itemize}
            \item $\tau \in D_{\varphi}$ and $ \varphi(\tau) = \xi $
        \end{itemize}

        Then we call it an initial value problem (Cauchy problem).
\end{description}
\section{Linear and higher-order linear differential equations}
\subsection{Linear differential equations}
\begin{description}
    \item[Definition] \hfill \\
        A linear differential equation is a differential equation for which:\\
        $ n=1, \quad I,I_1 \subset \R $ open intervals, $f:I\times I_1 \rightarrow \R$, where \\
        $g,h : I \rightarrow \R, \ g,h \in C, \ I_1 := \R$ and \\
        $f(x,y) := g(x)\cdot y + h(x) \quad (x \in I, y \in I_1 = \R) $\\
        $ \Rightarrow \varphi'(x) = f(x, \varphi(x)) = g(x) \cdot \varphi(x) + h(x) \quad (x \in D_{\varphi})$
    \item[Homogeneity] \hfill \\
        The linear differential equation is homogeneous if $ h \equiv 0$ (otherwise inhomogeneous)
    \item[Initial value problem] \hfill
        \begin{itemize}
            \item Every initial value problem relating to a linear differential equation is solvable and \\
                  for $\forall \varphi, \psi $ solutions: $ \varphi(t) = \psi(t) \quad (t \in D_{\varphi} \cap D_{\psi} )$
            \item Every solution of a homogeneous linear differential equation ($\varphi : I \rightarrow \R$) has the following form: \\
                  $ c\varphi_0$, where \\
                  $c \in \R$ and $\varphi_0(t) = e^{G(t)} \quad (G:I\rightarrow\R, \ G \in D, $ and $ G' = g)$
            \item Method of variation of constants:\\
                  $ \exists m:I\rightarrow\R, \ m \in D : m\cdot\varphi_0$ is a solution of the (inhomogeneous) linear differential equation
            \item Particular solution: \\
                  $M := \{ \varphi : I \rightarrow \R : \varphi'(t) = g(t)\cdot\varphi(t) + h(t) \ (t \in I)\} $ \\
                  $M_h := \{ \varphi : I \rightarrow \R : \varphi'(t) = g(t)\cdot\varphi(t) \ (t \in I)\} $\\
                  $\Rightarrow \forall \psi \in M : M = \psi + M_h = \{\varphi + \psi : \varphi \in M_h\}$\\
                  (And here $\psi$, based on the above, can be written in the form $m\cdot\varphi_0$)
            \item Example: Radioactive decay: \\
                  $ m_0 > 0$ - initial amount of material \\
                  $ m \in \R \rightarrow \R $ - mass-time function, where \\
                  $m(t)$ - the amount of material present \\
                  $ m \in D \Rightarrow \dfrac{m(t) - m(t+\Delta t)}{\Delta t} \quad (\Delta t \neq 0) $ - average decay rate \\
                  $ \dfrac{m(t) - m(t+\Delta t)}{\Delta t} \xrightarrow[\Delta t \rightarrow 0]{} -m'(t) $, which based on observation is $ \approx m(t)$ \\

                  that is: \\
                  $ m'(t) = - \alpha \cdot m(t) \quad (t\in\R, 0 < \alpha \in \R)$\\
                  $ m(0) = m_0 $ \\
                  \rule{3cm}{0.2pt} \\
                  Homogeneous linear differential equation (initial value problem): \\
                  $ g \equiv -\alpha, \ \tau :=0, \ \xi := m_0 $ \\
                  $ \Rightarrow G(t) = -\alpha t \quad (t\in\R) \Rightarrow \varphi_0(t) = e^{-\alpha t} \quad (t \in \R)$ \\
                  $ \Rightarrow \exists c \in \R : m(t) = c\cdot e^{-\alpha t} \quad (t \in \R)$, where \\
                  $m(0) = c = m_0 \Longrightarrow m(t) = m_0e^{-\alpha t} \quad (t \in \R)$ \\
                  If $ T \in \R : m(T) = \dfrac{m_0}{2} $ (half-life) \\
                  $\Rightarrow \dfrac{m_0}{2} = m_0e^{-\alpha T} \Rightarrow \dfrac{1}{2} = e^{-\alpha T} \Rightarrow e^{\alpha T} = 2$ \\
                  $\Rightarrow T = \dfrac{ln(2)}{\alpha} $
        \end{itemize}
\end{description}
\subsection{Higher-order linear differential equations}
\begin{description}
    \item[Definition] \hfill \\
        $ 0 < n \in \mathbb{N}, I \subset \R$ open, $ a_0, ... ,a_{n-1} :I \rightarrow \R$ continuous and $ c: I \rightarrow \R$ continuous. \\
        Let us look for a function $ \varphi \in I \rightarrow \mathbb{K}$ such that:
        \begin{itemize}
            \item $ \varphi \in D^n$
            \item $ D_{\varphi}$ open interval
            \item $ \varphi^{(n)}(x) + \sum\limits_{k=0}^{n-1}a_k(x) \cdot \varphi^{(k)}(x) = c(x) \quad (x \in D_{\varphi}) $
        \end{itemize}

        We call this an $n$-th order linear differential equation. (For $n=1$, a linear differential equation.) If additionally: \\
        $ \tau \in I, \ \xi_0, ... , \xi_{n-1} \in \mathbb{K}$ and
        \begin{itemize}
            \item $ \tau \in D_{\varphi}$ and $ \varphi^{(k)}(\tau) = \xi_k \quad (k = 0...n-1) $
        \end{itemize}
        Then we speak of an initial value problem.

    \item[Homogeneity] \hfill \\
        If $c(x) = 0$ we speak of a homogeneous $n$-th order linear differential equation. So the solution sets of the homogeneous and inhomogeneous equations: \\
        $ M_h := \{\varphi : I \rightarrow \mathbb{K} : \varphi \in D^n, \ \varphi^{(n)} + \sum\limits_{k=0}^{n-1}a_k\cdot\varphi^{(k)} = 0 \} $ \\
        $ M := \{\varphi : I \rightarrow \mathbb{K} : \varphi \in D^n, \ \varphi^{(n)} + \sum\limits_{k=0}^{n-1}a_k\cdot\varphi^{(k)} = c \} $ \\
        (Here $M_h$ is an $n$-dimensional linear space, so some $ \varphi_1,...,\varphi_n \in M_h$ form a basis, also called a fundamental system.)
    \item[Constant-coefficient case] \hfill \\
        In this case $a_0,...,a_{n-1} \in \R$
        \begin{itemize}
            \item Role of the characteristic polynomial \\
                  Let $P(t) := t^n + \sum\limits_{k=0}^{n-1}a_kt^k \quad (t \in \mathbb{K})$ be the characteristic polynomial and \\
                  $ \varphi_\lambda(x) := e^{\lambda x} \quad (x \in \R, \lambda \in \mathbb{K}) $ \\\\
                  Then: $ \varphi_\lambda \in M_h \Longleftrightarrow P(\lambda) = 0 $\\
                  Moreover, if $ \lambda $ is an $r$-fold root of $P$, and \\
                  $ \varphi_{\lambda,j}(x) := x^je^{\lambda x} \ (j = 0..r-1, x\in\R)$, then:
                  $ \varphi_{\lambda,j} \in M_h \Longleftrightarrow \varphi_{\lambda, j}^{(n)}+\sum\limits_{k=0}^{n-1}a_k\varphi_{\lambda, j}^{(k)} $ \\
                  that is $P(\lambda)^{(j)} = 0 \quad (j = 0..r-1)$
            \item Real solutions \\
                  Let $ \lambda = u+iv \quad (u,v \in \R, v\neq0, i^2 = -1) $ \\
                  $ \Rightarrow $ the functions $ x \mapsto x^je^{ux}cos(vx)$, and $x \mapsto x^je^{ux}sin(vx)$ form a real fundamental system (basis) (in $M_h$)
        \end{itemize}
    \item[Example: Oscillations] \hfill \\
        Describe the motion of a point mass of mass $m$ performing oscillatory motion along a straight line around a fixed point, if we know the position it occupied at the start of the observation and its velocity at that time! \\
        $ \varphi \in \R \rightarrow \R, \varphi \in D^2$ : displacement-time function \\
        $ m > 0 $ : mass \\
        $ F \in \R \rightarrow \R $ : displacing force \\
        $ \alpha > 0 $ : restoring force, proportional to $ \varphi $ \\
        $ \beta \geq 0 $ : damping force, proportional to the velocity. \\
        $ \Longrightarrow $ (based on Newton's law of motion):\\
        $ m \cdot \varphi'' = F - \alpha\varphi-\beta\varphi'$\\
        $ \varphi(0) = s_0, \varphi'(0) = s'_0 $\\
        \rule{4cm}{0.2pt} \\
        Second-order linear differential equation (initial value problem)\\
        Written in standard form: $ \varphi'' + \dfrac{\beta}{m}\varphi' + \dfrac{\alpha}{m}\varphi = \dfrac{F}{m} $

        Consider the periodic external constraint as a forced oscillation, when: \\
        $ \dfrac{F(x)}{m} = Asin(\omega x )  \quad [A>0$ (amplitude), $ \omega > 0$ (forcing frequency)] \\
            Then $ \omega_0 := \sqrt{\dfrac{\beta}{m}} $ - natural frequency\\
            and $\varphi''(x) + \omega_0^2\varphi(x) = Asin(\omega x) $ \\
            whose characteristic polynomial is: $ P(t) = t^2+\omega_0^2 \quad (t \in \R) $ \\
            Its solutions: $ \lambda = \pm \ \omega_0i $ \\

            We saw earlier that if $ \lambda = u+iv$ then the functions $ x \mapsto x^je^{ux}cos(vx)$, and $x \mapsto x^je^{ux}sin(vx)$ form a real fundamental system (basis) (in $M_h$). So $ \varphi(x) = c_1cos(\omega_0x) + c_2sin(\omega_0x) $ can be written in a form which, using a phase angle, can be rewritten as: $ d\cdot sin(\omega_0x+\delta) \quad (d = \sqrt{c_1^2+c_2^2}, \delta \in \R)$. So: \\
        $ M_h =  \{ d\cdot sin(\omega_0x+\delta)\}$

        Then we can already easily give a particular solution:
        \begin{itemize}
            \item for $\omega \neq \omega_0$ the particular solution: \\
                  $ x \rightarrow q\cdot sin(\omega x) $\\
                  And $ q = \dfrac{A}{\omega_0^2-\omega^2} $ satisfies the equation $-q\omega^2sin(\omega x)+\omega_0^2q\cdot sin(\omega x) = Asin(\omega x) $.
                  So: \\
                  $ \varphi(x) = d\cdot sin(\omega_0x + \delta)+\dfrac{A}{\omega_0^2-\omega^2}sin(\omega x) $ the solution is the sum of two harmonic oscillations.
            \item for $ \omega = \omega_0 $ (resonance) the particular solution: \\
                  $ x \rightarrow qx\cdot cos(\omega x) $\\
                  And $ q = \dfrac{-A}{2\omega} $ satisfies the equation $-2q\omega \cdot sin(\omega x)- q\omega^2x\cdot cos(\omega x) +\omega^2qx\cdot cos(\omega x) = Asin(\omega x) $.
                  So: \\
                  $ \varphi(x) = d\cdot sin(\omega x + \delta)-\dfrac{A}{2\omega}x\cdot cos(\omega x) $ the solution is the sum of a harmonic and an aperiodic oscillation.\\
                  (In this case, as time (x) passes, the value of $\varphi $ grows. In certain models this causes the ``disintegration of the system''.)
        \end{itemize}
\end{description}
\end{document}
